% Job posting: https://ca.linkedin.com/jobs/view/computer-vision-engineer-at-assembler-ai-4459590886
\documentclass[11pt]{article}
\usepackage[left=0.7in,top=0.7in,right=0.7in,bottom=0.7in]{geometry}
\usepackage{enumitem}
\usepackage[hidelinks]{hyperref}
\usepackage{titlesec}
\usepackage{array}
\usepackage{setspace}
\usepackage{needspace}
\usepackage[T1]{fontenc}
\usepackage{newtxtext,newtxmath}
\ifdefined\pdfgentounicode
  \input{glyphtounicode}
  \pdfgentounicode=1
\fi
\setstretch{1.04}
\pagenumbering{gobble}
\titleformat{\section}{\Large\scshape}{}{4em}{}[\vspace{-0.7em}\rule{\linewidth}{0.4pt}\vspace{-0.4em}]
\titlespacing*{\section}{0pt}{1em}{0.4em}
\newenvironment{cvrole}[5]{%
  \par\noindent\textbf{\large #1} | \textit{\small #2, #3}\hfill \textit{\small #4 - #5}\par
  \begin{itemize}[leftmargin=2em, itemsep=-0.3em, topsep=-0.5em]%
}{%
  \end{itemize}\par\noindent\mbox{}\par\vspace{0.1em}%
}
\newcommand{\cvName}{Nishal Stanislaus Silva}
\newcommand{\cvEmail}{nishal.silva@hotmail.com}
\newcommand{\cvPhone}{+1 613 200 2030}
\newcommand{\cvCity}{Toronto, ON}
\newcommand{\cvSite}{https://nishal.xyz}
\newcommand{\cvLinkedIn}{https://www.linkedin.com/in/nishal-silva}
\newcommand{\cvGitHub}{https://github.com/ns2max}
\newcommand{\cvStatus}{Canadian Permanent Resident, Eligible to work in Canada without sponsorship}
\newcommand{\cvTagline}{Computer Vision \& ML Engineer | Industrial Inspection, IoT \& Edge Deployment}
\begin{document}
\pagestyle{empty}
\begin{center}
    {\LARGE \scshape \textbf{\cvName}}{\large \textit{, PhD}}\\[1pt]
    {\small \cvTagline}\\[1pt]
    {\footnotesize\textit{\cvStatus}}\\[1pt]
    {\small
    \href{mailto:\cvEmail}{\cvEmail} \,|\, \cvPhone \,|\, \cvCity
    \,|\, \href{\cvSite}{nishal.xyz} \,|\, \href{\cvLinkedIn}{linkedin.com/in/nishal-silva}
    \,|\, \href{\cvGitHub}{github.com/ns2max}}
\end{center}

\section*{Professional Summary}
Computer vision engineer with 5.5 years building and deploying production machine-vision systems on live manufacturing floors at MAS Holdings, South Asia's largest apparel manufacturer -- training and evaluating vision models end to end, from raw camera footage through annotation, evaluation-framework design and deployment, on real production data rather than curated benchmarks. Automated a multilingual care-label QC system with OpenCV template and feature matching driving a PLC-controlled conveyor reject, cutting inspection time 99.5\%, and a microscopic camera-array system for loom-side fabric defect detection that cut operator headcount 50\%. A PhD in Information Engineering and Computer Science (University of Trento) added rigor in evaluation methodology, benchmark and ablation design, and real-time embedded inference: a published detector running at 14 ms on a Raspberry Pi 4, F1 0.76, 74x faster than a DTW baseline. Strong Python and PyTorch proficiency, hands-on OpenCV and ARM/embedded deployment experience under hard latency budgets, and a track record of investigating model failures against real-world footage rather than lab conditions. Canadian Permanent Resident, no sponsorship required, based in Toronto and open to relocating to Montreal.

\section*{Core Competencies}
\textbf{Computer Vision \& Perception:} training and evaluating vision models on production video and camera footage (PyTorch, OpenCV), template and feature matching, image registration and warping, corner/edge detection, explainable defect overlays for human-in-the-loop QC, evaluation-framework design for real deployment conditions (precision/recall trade-offs, train/test protocol design). \\[2pt]
\textbf{Data \& Model Lifecycle:} owning training datasets end to end, from sourcing and annotation-quality review through failure-mode investigation on live footage; running end-to-end experiments; industrial and microscopic camera-array systems; multilingual label and print inspection. \\[2pt]
\textbf{Embedded \& Real-Time Systems:} Raspberry Pi and ARM CPU profiling, C++17 real-time inference engines, hard latency-budget deployment (14 ms on-device), PLC integration for automated reject/sort actuation, high-frequency IoT sensor ingestion and ETL, Linux/SSH-based development.

\section*{Technical Stack}
Python, PyTorch, OpenCV, C++17, TensorFlow, Keras, scikit-learn, NumPy, pandas, Raspberry Pi / ARM, AWS, Git, Docker, CI/CD, pytest.

\section*{Education}
\noindent\textbf{PhD, Information Engineering and Computer Science} -- University of Trento, Italy \hfill 2021 - 2025 \\
\noindent\textbf{MSc, Telecommunication and Electronic Engineering} -- Sheffield Hallam University, UK \hfill 2016 - 2018 \\
\noindent\textbf{BEng (Hons), Electronic Engineering} -- Sheffield Hallam University, UK \hfill 2013 - 2014

\section*{Research and Work Experience}

\begin{cvrole}{Independent Researcher}{Self-directed}{Toronto, Canada}{Jan 2026}{Present}
\item Two peer-reviewed papers accepted (Smart Drums, JAES; MIRaaS, IEEE IS2 2026) covering real-time inference and edge deployment.
\item Maintain Nebula, an open-source C++17 audio-feature-extraction library with per-feature ARM/Raspberry Pi latency benchmarks; used AI-assisted development workflows, including Claude Code, throughout.
\end{cvrole}

\begin{cvrole}{Postdoctoral Researcher}{University of Trento}{Trento, Italy}{Jan 2025}{Dec 2025}
\item Owned the full ML pipeline for streaming audio and sensor data: acquisition, feature engineering, training and evaluation, edge deployment on Raspberry Pi 4, and monitoring.
\item Published a real-time detector achieving 14 ms inference on Raspberry Pi 4, F1 0.76, 31.4\% CPU utilization -- 74x faster than a DTW baseline (JAES 2026).
\item Designed frozen-protocol benchmark suites and ablation studies quantifying accuracy/latency/CPU trade-offs for embedded deployment.
\end{cvrole}

\begin{cvrole}{Visiting Researcher}{McGill University, IDMIL/CIRMMT}{Montreal, Canada}{May 2024}{Aug 2024}
\item Built a self-powered embedded sensor system (BNO055 IMU + HiFiBerry + Raspberry Pi 4); fused IMU and audio signals via a hierarchical finite-state machine, achieving F1 0.78 on a 500-event corpus (IEEE IS2 2025).
\end{cvrole}

\begin{cvrole}{Visiting Researcher}{University of Visual and Performing Arts}{Colombo, Sri Lanka}{Jan 2024}{Mar 2024}
\item Conducted a human-centred evaluation of a real-time detection system against user footage and feedback (IEEE I3DA 2025).
\end{cvrole}

\begin{cvrole}{Technical Consultant}{Forestpin (Pvt) Ltd}{Colombo, Sri Lanka}{Aug 2020}{Feb 2021}
\item Built ML and statistical pipelines for financial forensics, compliance monitoring, and anomaly detection on large structured enterprise data.
\end{cvrole}

\begin{cvrole}{Research Engineer}{MAS Holdings (Pvt) Ltd}{Colombo, Sri Lanka}{Jan 2015}{Jul 2020}
\item Designed and deployed end-to-end computer-vision and IoT systems for manufacturing at scale, taking camera and sensor streams directly into production workflows; drove operational efficiency gains of up to 300\%.
\item Built an automated multilingual care-label QC system using OpenCV template and feature matching with an explainable defect overlay; iterated from handheld scanner to industrial camera to a conveyor with PLC-controlled reject, cutting inspection time 99.5\%, with a human-in-the-loop path for borderline cases and continuous failure-mode review against production footage.
\item Developed a loom-side fabric structural-defect detection system using a microscopic camera array, cutting operator headcount 50\%.
\item Owned training data end to end for these systems: sourcing, annotation review, and iterative retraining against real production footage rather than a fixed benchmark set.
\item Introduced CI/CD and code review practices; mentored engineers and delivered group-wide computer-vision training.
\item Designed and deployed a remote patient-vitals computer-vision monitoring device for COVID-19 hospital wards, recognized by Sri Lanka's Government Medical Officers' Association.
\end{cvrole}

\section*{Publications}
Silva, N. \& Turchet, L. (2026). Real-Time Audio Pattern Detection for Smart Musical Instruments. \textit{Journal of the Audio Engineering Society}, 74(3). Silva, N. \& Turchet, L. (2026, in press). Smart Drums. \textit{Journal of the Audio Engineering Society}.

\section*{Languages}
English (native/C2) \,|\, Sinhala (mother tongue) \,|\, Italian (B1) \,|\, French (basic, A2, actively learning)

\end{document}
